% SGA 3, Expose V, Loop 1 English reconstruction.
% Scope: component-local PDF pp. 6--10, whose page headers are printed
% pp. 256--260. This begins with the continuation of Section 2(a) and ends
% at the end of Section 3(f), including editor's notes 9--15. The legacy
% marginal page markers advance through 262 inside this five-page range.
% Component-local p. 11 (Section 4 / Theorem 4.1) is excluded.
% Authority: Polo--Gille Exp5-13oct24.pdf, combined-reader pp. 268--272.
% OCR and Jacob Reinhold's English Markdown are comparison witnesses only.

% Source page: Exp5 component-local p. 6; combined-reader p. 268;
% page header p. 256.
\(G\) 的 \(\mathscr C\)-群结构的箭头, 最后用
\(\operatorname{pr}_{2,3}\) 表示
\(G\times G\times X=G\times(G\times X)\) 到第二个因子的投影. 则

\SGAThreeDiagram[0.47\textwidth]
  {figures/exp5/p256_action_groupoid_continuation.png}
  {与 \(\mathscr C\)-群 \(G\) 在 \(X\) 上的作用关联的群胚.}
  {fig:sga3-exp5-p256-action-groupoid}
  {Polo--Gille Exp5-13oct24.pdf, 组件第 6 页; 合卷第 268 页; 页眉第 256 页}

\noindent 是一个 \(\mathscr C\)-群胚.

\medskip
\phantomsection\label{par:sga3-V-2-b}%
\noindent\textbf{b)} 设 \(d_0,d_1:X_1\to X_0\) 是一个\emph{等价对}.
换言之, 若 \(d_0\boxtimes d_1:X_1\to X_0\times X_0\)
是分量为 \(d_0\) 和 \(d_1\) 的箭头, 则假设对
\(\mathscr C\) 的每个对象 \(T\), 映射 \((d_0\boxtimes d_1)(T)\)
是从 \(X_1(T)\) 到 \(X_0(T)\) 上某个等价关系的图像的双射.
因此, 将 \(X_1(T)\) 等同于 \(X_0(T)\) 中满足 \(x\sim y\)
的元素对 \((x,y)\) 的集合. 同样地,
\[
  X_2(T)=\bigl(X_1\times_{d_1,d_0}X_1\bigr)(T)
\]
可等同于 \(X_0(T)\) 中满足 \(x\sim y\) 和 \(y\sim z\)
的元素三元组 \((x,y,z)\) 的集合. 因而有唯一箭头
\(d'_1:X_2\to X_1\), 使方形
\hyperref[sga3:v:lemma:1-2:equation:1]{(1)} 和
\hyperref[sga3:v:lemma:1-2:equation:2]{(2)} 交换:
\(d'_1(T)\) 必须将 \((x,y,z)\in X_2(T)\) 映到 \((x,z)\in X_1(T)\).
作此选择的 \(d'_1\) 使
\[
  \bigl(d_0,d_1:X_1\rightrightarrows X_0,d'_1\bigr)
\]
成为 \(\mathscr C\)-群胚.

反之, 考虑一个满足 \(d_0\boxtimes d_1:X_1\to X_0\times X_0\)
为单态射的 \(\mathscr C\)-群胚 \(X_\ast\).
则 \((d_0,d_1)\) 是等价对, 并且可按
\hyperref[par:sga3-V-2-b]{上面几行}所述, 从 \((d_0,d_1)\)
重建 \(X_\ast\).%
\footnote{\label{note:sga3-V-9}编者注: 特别地, 若 \(G\)
是左作用在 \(\mathscr C\) 的对象 \(X\) 上的 \(\mathscr C\)-群,
\(X_\ast\) 是 \hyperref[par:sga3-V-2-a]{2(a)} 中定义的
\(\mathscr C\)-群胚, 则 \((d_0,d_1)\) 为等价对当且仅当
\(G\) 在 \(X\) 上自由作用; 参见第 III 篇, 3.2.1.}

\medskip
\phantomsection\label{par:sga3-V-2-c}%
\noindent\textbf{c)} 若 \(p:X\to Y\) 是 \(\mathscr C\) 的任意箭头,
\(\operatorname{pr}_1\) 和 \(\operatorname{pr}_2\) 是
\(X\times_{p,p}X\) 到 \(X\) 的两个投影, 则
\[
  (\operatorname{pr}_1,\operatorname{pr}_2):
  X\times_{p,p}X\rightrightarrows X
\]
是一个等价对. 若图

\SGAThreeDiagram[0.36\textwidth]
  {figures/exp5/p256_effective_epimorphism.png}
  {定义有效满态射的图.}
  {fig:sga3-exp5-p256-effective-epimorphism}
  {Polo--Gille Exp5-13oct24.pdf, 组件第 6 页; 合卷第 268 页; 页眉第 256 页}

\noindent 正合, 即
\[
  (Y,p)=\operatorname{Coker}(\operatorname{pr}_1,\operatorname{pr}_2).
\]
则称箭头 \(p\) 为\emph{有效满态射}.

例如, 设 \(S\) 是 Noether 概形, \(\mathscr C\) 是
\(S\) 上有限的概形构成的范畴. 我们说明 \(\mathscr C\)
中的满态射未必有效. 取
\[
  S=\operatorname{Spec}k[T^3,T^5],\qquad
  Y=S,\qquad X=\operatorname{Spec}k[T],
\]
其中 \(k\) 是交换域. 若 \(i\) 是 \(B=k[T^3,T^5]\)
到 \(A=k[T]\) 的包含映射, 则取 \(p=\operatorname{Spec}i\).
于是 \(X\times_{p,p}X\) 可等同于 \(\operatorname{Spec}(A\otimes_B A)\),
\(\operatorname{Coker}(\operatorname{pr}_1,\operatorname{pr}_2)\)
可等同于 \(\operatorname{Spec}B'\), 其中 \(B'\) 是 \(A\)
中由满足 \(a\otimes_B1=1\otimes_Ba\) 的元素 \(a\) 组成的子环. 但
\[
  T^7\otimes_B1
  =(T^2T^5)\otimes_B1
  =T^2\otimes_BT^5
  =T^2\otimes_B(T^3T^2)
  =T^5\otimes_BT^2
  =1\otimes_BT^7.
\]
因此, \(T^7\) 属于 \(B'\), 不属于 \(B\), 且
\(\operatorname{Spec}B'\) 不同于 \(\operatorname{Spec}B\),
这就给出了反例.%
\footnote{\label{note:sga3-V-10}编者注: 相同论证适用于
\(B=k[T^3,T^4]\) 和 \(T^5\otimes_B1\);
更一般地, 它适用于 \(B=k[T^n,T^{n+r}]\) 和元素
\(T^{n+2r}\otimes_B1\), 只要 \(n\) 不整除 \(2r\).}

% Source page: Exp5 component-local p. 7; combined-reader p. 269;
% page header p. 257.
\section*{3. 关于 \(\mathscr C\)-群胚的若干杂论}
\phantomsection\label{sec:sga3-V-3}
\addcontentsline{toc}{section}{3. 关于 \texorpdfstring{\(\mathscr C\)}{C}-群胚的若干杂论}

下面不按特定次序列出后文要用到的一些注记.

\medskip
\phantomsection\label{par:sga3-V-3-a}%
\noindent\textbf{a)} 设

\SGAThreeDiagram[0.31\textwidth]
  {figures/exp5/p257_x_groupoid.png}
  {\(\mathscr C\)-群胚 \(X_\ast\).}
  {fig:sga3-exp5-p257-x-groupoid}
  {Polo--Gille Exp5-13oct24.pdf, 组件第 7 页; 合卷第 269 页; 页眉第 257 页}

\noindent 是一个 \(\mathscr C\)-群胚, \(f_0:Y_0\to X_0\)
是 \(\mathscr C\) 的箭头. 我们将定义以 \(Y_0\) 为底的
\(\mathscr C\)-群胚

\SGAThreeDiagram[0.31\textwidth]
  {figures/exp5/p257_y_induced_groupoid.png}
  {诱导的 \(\mathscr C\)-群胚 \(Y_\ast\).}
  {fig:sga3-exp5-p257-y-induced-groupoid}
  {Polo--Gille Exp5-13oct24.pdf, 组件第 7 页; 合卷第 269 页; 页眉第 257 页}

\noindent 称之为\emph{由 \(X_\ast\) 和 \(f_0\) 诱导的群胚}.
也称 \(Y_\ast\) 为 \emph{\(X_\ast\) 在基变换 \(f_0\) 下的逆像}.

选取 \(Y_1\) 为下图的纤维积

\SGAThreeDiagram[0.29\textwidth]
  {figures/exp5/p257_base_change_fiber_square.png}
  {\(Y_1\) 的纤维积定义.}
  {fig:sga3-exp5-p257-base-change-fiber-square}
  {Polo--Gille Exp5-13oct24.pdf, 组件第 7 页; 合卷第 269 页; 页眉第 257 页}

\noindent 并取 \(e_0\) 和 \(e_1\) 分别为典范箭头
\(Y_1\to Y_0\times Y_0\) 与 \(Y_0\times Y_0\)
的第一个和第二个投影的复合. 于是态射 \(Y_1\to Y_0\times Y_0\)
就是 \(e_0\boxtimes e_1\), 且
\[
  f_0\circ e_i=d_i\circ f_1\qquad(i=0,1),
\]
其中 \(f_1\) 表示 \(Y_1\) 到 \(X_1\) 的投影.

令 \(Y_2=Y_1\times_{e_0,e_1}Y_1\), 参见
\hyperref[par:sga3-V-1-c]{1(c)}.
可以说, 箭头对 \((e_0,e_1)\) 的定义使得对每个
\(T\in\mathscr C\) 和 \(Y_0(T)\) 的每个元素对 \((y,x)\),
在 \(X_\ast(T)\) 中源为 \(f_0(x)\)、目标为 \(f_0(y)\) 的箭头 \(\psi\),
与 \(Y_\ast(T)\) 中源为 \(x\)、目标为 \(y\) 的箭头 \({}_{y}\psi_x\)
之间, 有一个特定的双射对应
\[
  \psi\longmapsto{}_{y}\psi_{x}
\]
因此, 对每个 \(T\in\mathscr C\), 用公式
\[
  {}_{z}\psi_y\circ{}_{y}\varphi_x
  ={}_{z}(\psi\circ\varphi)_x.
\]
定义 \(Y_\ast(T)\) 中箭头的复合, 就确定了 \(e'_1:Y_2\to Y_1\).
显然, 这一定义使每个 \(Y_\ast(T)\) 都成为群胚.

\medskip
\phantomsection\label{par:sga3-V-3-b}%
\noindent\textbf{b)} 给定 \(\mathscr C\)-群胚 \(X_\ast\)
和基变换 \(f_0:Y_0\to X_0\), 还可用另一种方式重建箭头对
\((e_0,e_1):Y_1\rightrightarrows Y_0\).%
\footnote{\label{note:sga3-V-11}编者注: 这一观点用于
\hyperref[par:sga3-V-3-f]{3(f)} 及 \hyperref[lem:sga3-V-6.1]{6.1} 的证明.}
构造 \(Y_0\times_{X_0}X_1\), \(\operatorname{pr}_1\)
和 \(\operatorname{pr}_2\), 使方形

\SGAThreeDiagram[0.27\textwidth]
  {figures/exp5/p257_reconstruction_square.png}
  {定义 \(Y_0\times_{X_0}X_1\) 的 笛卡尔 方形.}
  {fig:sga3-exp5-p257-reconstruction-square}
  {Polo--Gille Exp5-13oct24.pdf, 组件第 7 页; 合卷第 269 页; 页眉第 257 页}

% Source page: Exp5 component-local p. 8; combined-reader p. 270;
% page header p. 258.
\noindent 为 笛卡尔 方形. 直接化归到集合的情形后,
可知下面的方形也是 笛卡尔 方形:

\SGAThreeDiagram[0.27\textwidth]
  {figures/exp5/p258_cartesian_reconstruction_square.png}
  {重建 \((e_0,e_1)\) 时的第二个 笛卡尔 方形.}
  {fig:sga3-exp5-p258-cartesian-reconstruction-square}
  {Polo--Gille Exp5-13oct24.pdf, 组件第 8 页; 合卷第 270 页; 页眉第 258 页}

这里 \(f_1\) 表示
\[
  Y_1=(Y_0\times Y_0)\times_{(X_0\times X_0)}X_1
\]
到 \(X_1\) 的典范投影.

\medskip
\phantomsection\label{par:sga3-V-3-c}%
\noindent\textbf{c)} 给出 \(\mathscr C\)-群胚的逆像的两个例子.
取 \(Y_0=X_1\), \(f_0=d_0\).
对 \(\mathscr C\) 的每个对象 \(T\), 集合 \(Y_1(T)\)
可等同于 \(X_\ast(T)\) 中形如

\SGAThreeDiagram[0.18\textwidth]
  {figures/exp5/p258_inverse_image_phi_diagram.png}
  {第一个逆像群胚的箭头.}
  {fig:sga3-exp5-p258-inverse-image-phi}
  {Polo--Gille Exp5-13oct24.pdf, 组件第 8 页; 合卷第 270 页; 页眉第 258 页}

\noindent 的图的集合. 这样的图的源是箭头 \(f\), 目标是箭头 \(g\).
这些图以显然的方式复合.

现令 \(Y'_0=X_1\), \(f'_0=d_1\)
(添加撇号是为了避免与前一个例子混淆).% 
\footnote{\label{note:sga3-V-12}编者注: 原文用语为``附加记号''.}
此时, 对每个 \(T\in\mathscr C\), 集合 \(Y'_1(T)\)
可等同于群胚 \(X_\ast(T)\) 中形如

\SGAThreeDiagram[0.18\textwidth]
  {figures/exp5/p258_inverse_image_psi_diagram.png}
  {第二个逆像群胚的箭头.}
  {fig:sga3-exp5-p258-inverse-image-psi}
  {Polo--Gille Exp5-13oct24.pdf, 组件第 8 页; 合卷第 270 页; 页眉第 258 页}

\noindent 的图的集合. 这样的图的源为 \(f\), 目标为 \(g\),
这些图同样以显然的方式复合.

现在显然, \(Y_0(T)\) 的恒等映射与映射

\SGAThreeDiagram[0.44\textwidth]
  {figures/exp5/p258_groupoid_isomorphism_map.png}
  {从 \(Y_1(T)\) 到 \(Y'_1(T)\) 的映射.}
  {fig:sga3-exp5-p258-groupoid-isomorphism-map}
  {Polo--Gille Exp5-13oct24.pdf, 组件第 8 页; 合卷第 270 页; 页眉第 258 页}

\noindent 一起定义从群胚 \(Y_\ast(T)\) 到 \(Y'_\ast(T)\) 的同构.
此外, 此同构关于 \(T\) 具有函子性,
所以 \(\mathscr C\)-群胚 \(Y_\ast\) 与 \(Y'_\ast\) 同构.%
\footnote{\label{note:sga3-V-13}编者注:
这在 \hyperref[lem:sga3-V-6.1]{引理~6.1} 的证明中起关键作用.}

% Source page: Exp5 component-local p. 9; combined-reader p. 271;
% page header p. 259.
\medskip
\phantomsection\label{par:sga3-V-3-d}%
\noindent\textbf{d)}

\medskip
\phantomsection\label{prop:sga3-V-3.1}%
\noindent\textbf{命题 3.1.}---
\textit{沿用 \hyperref[par:sga3-V-3-a]{3(a)} 的记号,
并假设 \(f_0\) 是既有效又普遍的满态射.
则 \(\operatorname{Coker}(d_0,d_1)\) 存在当且仅当
\(\operatorname{Coker}(e_0,e_1)\) 存在.%
\footnote{\label{note:sga3-V-14}编者注: 修改了原文, 以明确写出下面的同构.}
此外, 在这种情形下, \(f_0\) 诱导一个同构}
\phantomsection\label{eq:sga3-V-3.1-cokernel-isomorphism}%
\[
  \operatorname{Coker}(d_0,d_1)
  \xrightarrow{\ \sim\ }
  \operatorname{Coker}(e_0,e_1).
\]

先回顾一下, 满态射 \(f_0:Y_0\to X_0\) 称为\emph{普遍的},
是指在每个 笛卡尔 方形

\SGAThreeDiagram[0.18\textwidth]
  {figures/exp5/p259_universal_epimorphism_square.png}
  {普遍满态射 \(f_0\) 的一个基变换.}
  {fig:sga3-exp5-p259-universal-epimorphism-square}
  {Polo--Gille Exp5-13oct24.pdf, 组件第 9 页; 合卷第 271 页; 页眉第 259 页}

\noindent 中, 箭头 \(f'\) 都是满态射.
现用 \(C(d_0,d_1)\) 表示从 \(\mathscr C\) 到集合范畴的协变函子,
它对每个 \(T\in\mathscr C\) 赋予箭头对
\[
  T(d_0),T(d_1):T(X_0)\rightrightarrows T(X_1),
\]
的核, 并类似地定义 \(C(e_0,e_1)\).
因而对每个 \(T\in\mathscr C\), 有交换图

\SGAThreeDiagram[0.51\textwidth]
  {figures/exp5/p259_functor_kernel_diagram.png}
  {函子 \(C(d_0,d_1)\) 与 \(C(e_0,e_1)\) 的比较.}
  {fig:sga3-exp5-p259-functor-kernel-diagram}
  {Polo--Gille Exp5-13oct24.pdf, 组件第 9 页; 合卷第 271 页; 页眉第 259 页}

这里 \(T(f)\) 是单射 \(T(f_0)\) 诱导的单射.
若证明 \(T(f)\) 对每个 \(T\) 都满射, 则得到一个函子性的同构
\[
  f:C(d_0,d_1)\xrightarrow{\ \sim\ }C(e_0,e_1).
\]
因此, 这两个函子中的一个可表当且仅当另一个可表,
这就证明了 \hyperref[prop:sga3-V-3.1]{本命题}.

为证明 \(T(f)\) 满射, 考虑下图

\SGAThreeDiagram[0.52\textwidth]
  {figures/exp5/p259_delta_proof_diagram.png}
  {用来证明 \(T(f)\) 的满射性的图.}
  {fig:sga3-exp5-p259-delta-proof-diagram}
  {Polo--Gille Exp5-13oct24.pdf, 组件第 9 页; 合卷第 271 页; 页眉第 259 页}

这里 \(\Delta\) 是由态射
\[
  s\circ f_0\circ\operatorname{pr}_1:Y_0\times Y_0\longrightarrow X_1,
\]
定义的 \(Y_1\to Y_0\times Y_0\) 的截面,
其中 \(s:X_0\to X_1\) 满足 \hyperref[sec:sga3-V-1]{第~1~节} 的
\hyperref[eq:sga3-exp5-3]{(3)} 和
\hyperref[eq:sga3-exp5-3bis]{(3 bis)}.

% Source page: Exp5 component-local p. 10; combined-reader p. 272;
% page header p. 260. Legacy marginal markers on this page reach 262.
若 \(g:Y_0\to T\) 满足 \(g\circ e_0=g\circ e_1\), 则
\[
  g\circ e_0\circ\Delta=g\circ e_1\circ\Delta,
\]
所以 \(g\circ\operatorname{pr}_1=g\circ\operatorname{pr}_2\).
由于 \(f_0\) 是有效满态射, \(g\) 可分解为 \(h\circ f_0\),
其中 \(h:X_0\to T\) 为某个箭头; 换言之, \(g=T(f_0)(h)\).
只须证明 \(h\in C(d_0,d_1)(T)\), 或等价地, \(hd_0=hd_1\).
但
\[
  hd_0f_1=hf_0e_0=ge_0=ge_1=hf_0e_1=hd_1f_1,
\]
所求等式由 \(f_1\) 是满态射得到 (因为 \(f_0\) 是普遍满态射).

\medskip
\phantomsection\label{par:sga3-V-3-e}%
\noindent\textbf{e)} 现设 \(S\) 为概形, 并取 \(\mathscr C=(\mathbf{Sch}/S)\).
给定 \(\mathscr C\)-群胚

\SGAThreeDiagram[0.37\textwidth]
  {figures/exp5/p260_scheme_groupoid.png}
  {一个 \((\mathbf{Sch}/S)\)-群胚.}
  {fig:sga3-exp5-p260-scheme-groupoid}
  {Polo--Gille Exp5-13oct24.pdf, 组件第 10 页; 合卷第 272 页; 页眉第 260 页}

\noindent 就在概形 \(X_0\) 的底集 \(\underline{X_0}\) 上定义了等价关系:
对 \(x,y\in\underline{X_0}\), 若存在 \(z\in\underline{X_1}\)
使 \(x=d_1z\) 及 \(y=d_0z\), 则记 \(x\sim y\).
自反性和对称性是显然的.%
\footnote{\label{note:sga3-V-15}编者注: 自反性来自箭头 \(s:X_0\to X_1\)
的存在性, 它同时为 \(d_0\) 和 \(d_1\) 的截面;
对称性来自 \(X_1\) 上``交换 \(d_0\) 与 \(d_1\)''的对合 \(\sigma\)
的存在性, 即满足 \(d_0\sigma=d_1\) 与 \(d_1\sigma=d_0\) 的对合;
参见 \hyperref[sec:sga3-V-1]{第~1~节},
\hyperref[eq:sga3-exp5-3]{(3)},
\hyperref[eq:sga3-exp5-3bis]{(3 bis)} 和
\hyperref[eq:sga3-exp5-dagger]{(\(\dagger\))}.}
证明传递性. 若 \(x\sim y\) 且 \(y\sim z\), 则存在
\(u,v\in\underline{X_1}\) 使
\[
  x=d_1u,\qquad y=d_0u,\qquad y=d_1v,\qquad z=d_0v.
\]
于是 \((v,u)\) 属于集合的纤维积
\(\underline{X_1}\times_{\underline{d_1},\underline{d_0}}
\underline{X_1}\). 由于典范映射

\SGAThreeDiagram[0.32\textwidth]
  {figures/exp5/p260_underlying_fiber_product_map.png}
  {从纤维积的底集到各底集的纤维积的典范映射.}
  {fig:sga3-exp5-p260-underlying-fiber-product-map}
  {Polo--Gille Exp5-13oct24.pdf, 组件第 10 页; 合卷第 272 页; 页眉第 260 页}

\noindent 满射, \((v,u)\) 是某个 \(w\in\underline{X_2}\) 的像.
因此
\[
  x=d_1d'_1w,\qquad z=d_0d'_1w,
\]
从而 \(x\sim z\).

\medskip
\phantomsection\label{par:sga3-V-3-f}%
\noindent\textbf{f)} 沿用 \hyperref[par:sga3-V-3-a]{3(a)} 和
\hyperref[par:sga3-V-3-b]{3(b)} 的记号, 仍取
\(\mathscr C=(\mathbf{Sch}/S)\). 若 \(x,y\) 为 \(Y_0\) 的点, 则
\[
  x\sim y\quad\Longleftrightarrow\quad f_0(x)\sim f_0(y).
\]
因此, 群胚所定义的等价关系的逆像,
就是该群胚的逆像所定义的等价关系.

事实上, 假设 \(x\sim y\). 存在 \(z\in\underline{Y_1}\)
使 \(x=e_1(z)\) 及 \(y=e_0(z)\).
由于对 \(i=0,1\) 有 \(f_0\circ e_i=d_i\circ f_1\), 可得
\[
  f_0(x)=d_1f_1(z),\qquad f_0(y)=d_0f_1(z),
\]
从而 \(f_0(x)\sim f_0(y)\).

反之, 假设 \(f_0(x)\sim f_0(y)\), 并令
\(z\in\underline{X_1}\) 满足
\[
  f_0(y)=d_1(z),\qquad f_0(x)=d_0(z).
\]
使用 \hyperref[par:sga3-V-3-b]{3(b)} 的构造和记号,
存在点 \(t\in\underline{Y_0\times_{X_0}X_1}\),
使 \(\operatorname{pr}_1(t)=x\) 及 \(\operatorname{pr}_2(t)=z\).
由于 \(f_0(y)=d_1\operatorname{pr}_2(t)\), 同样存在
\(s\in\underline{Y_1}\), 使
\[
  y=e_1(s),\qquad (e_0\boxtimes f_1)(s)=t.
\]
于是
\[
  e_0(s)=\operatorname{pr}_1(e_0\boxtimes f_1)(s)
  =\operatorname{pr}_1(t)=x,
\]
所以 \(x\sim y\).
% Translation statistics: source_lines=436; output_lines=426; inline_math=243; display_math=23; equation=0; equation_star=0; align=0; footnotes=7.
% Review: complete assigned component; pending_TODO=0; all math symbols, labels, references, diagram paths and environments checked; formula prose translated.
